\documentclass[11pt]{article}

\usepackage[margin=1in]{geometry}
\usepackage[T1]{fontenc}
\usepackage{amsmath,amssymb}
\usepackage{booktabs}
\usepackage{enumitem}
\usepackage{xcolor}
\usepackage{hyperref}

\hypersetup{
  colorlinks=true,
  linkcolor=blue!60!black,
  citecolor=blue!60!black,
  urlcolor=blue!60!black
}

\setlist[itemize]{leftmargin=1.5em}
\setlist[enumerate]{leftmargin=1.7em}

\title{Plan for Full-Ring Ultrasound Speed-of-Sound Reconstruction\\
Using a Tensorized Fourier Neural Operator}
\author{}
\date{September 2026}

\begin{document}

\maketitle

\begin{abstract}
This document describes a staged plan for training a Tensorized Fourier
Neural Operator (TFNO) to estimate a two-dimensional speed-of-sound field from
a complete circular ultrasound acquisition. Each example contains a known
transmit pulse, measurements from 256 transmitters and 256 receivers, and a
$512\times512$ target speed-of-sound image. The plan covers acquisition
geometry, full-wave simulation, storage, signal representation, network
architecture, training, validation, and software organization. Only the
recommended full-ring acquisition is considered.
\end{abstract}

\section{Objective and Scope}

The objective is to learn the nonlinear inverse operator
\begin{equation}
  \mathcal{G}^{-1}:\{s(t),d(\theta_{\mathrm{tx}},
  \theta_{\mathrm{rx}},t)\}\longmapsto c(x,y),
\end{equation}
where $s(t)$ is the transmitted pulse, $d$ is the measured pressure, and
$c(x,y)$ is the spatial speed of sound. The acquisition contains 256 transmit
events, with all 256 elements recording each event. Therefore, one unprocessed
measurement has the form
\begin{equation}
  d\in\mathbb{R}^{256\times256\times T},
\end{equation}
and the reconstruction target is
\begin{equation}
  c\in\mathbb{R}^{512\times512}.
\end{equation}

The proposed method follows the neural-operator formulation demonstrated for
ultrasound tomography by Dai et al.~\cite{dai2023}, while extending it from a
time-of-flight field to multi-frequency full-wave information and from 128 to
256 ring elements. Fourier neural operators learn maps between functions by
parameterizing integral kernels in Fourier space~\cite{li2021,kovachki2023}.
Tensorization reduces the number of spectral parameters and is available in
the official NeuralOperator implementation~\cite{kossaifi2024}.

\section{Fixed Acquisition Geometry}

All training, validation, test, and inference examples shall use one fixed
physical geometry. This is essential because changing the ring radius or grid
spacing changes the inverse operator that the network must learn.

The initial geometry will use:
\begin{itemize}
  \item a $512\times512$ Cartesian computational grid;
  \item isotropic spacing of approximately $1\,\mathrm{mm}$;
  \item a fixed ring radius of approximately $230\,\mathrm{mm}$;
  \item 256 equally spaced transmit-receive elements;
  \item water coupling outside the body;
  \item fixed element numbering and angular origin;
  \item a fixed recording interval and number of temporal samples.
\end{itemize}

Each CT-derived medium will be centered and padded into this common field of
view. It may be downsampled when necessary, but it will not be independently
rescaled to a different ring geometry. An image that cannot fit inside the
fixed ring will be rejected or handled by a single documented preprocessing
rule that is applied to the entire dataset.

\section{Full-Ring Forward Acquisition}

A new full-ring acquisition routine will wrap the existing
\texttt{simulate\_shot} function. For transmitter index $i$, the simulator
returns
\begin{equation}
  d_i(j,n),\qquad j=0,\ldots,255,
  \quad n=0,\ldots,T-1.
\end{equation}
Repeating this for all transmitters produces the complete tensor
\begin{equation}
  d(i,j,n)\in\mathbb{R}^{256\times256\times T}.
\end{equation}

The acquisition implementation will:
\begin{itemize}
  \item reset the wave solver before every transmit event;
  \item retain the transmitter self-channel but provide a validity mask;
  \item save transmitter and receiver coordinates, $\Delta t$, grid spacing,
        ring radius, pulse, and simulation configuration;
  \item support restart at the transmitter level after interruption;
  \item write samples incrementally rather than holding an entire dataset in
        memory;
  \item test acoustic reciprocity,
        $d_{ij}(t)\approx d_{ji}(t)$, as a simulation quality check.
\end{itemize}

Independent subjects or CT slices can be generated in parallel. Writes to a
single sample should remain serialized to avoid corrupting the dataset.

\section{Dataset Contract and Storage}

Each dataset sample will contain the following arrays and metadata:
\begin{center}
\begin{tabular}{lll}
\toprule
Field & Shape & Description \\
\midrule
\texttt{pulse} & $(T)$ & transmitted pressure waveform \\
\texttt{traces} & $(256,256,T)$ & full transmitter-receiver data \\
\texttt{sound\_speed} & $(512,512)$ & target $c(x,y)$ in $\mathrm{m/s}$ \\
\texttt{slowness\_squared} & $(512,512)$ & target $m=1/c^2$ \\
\texttt{ring\_mask} & $(512,512)$ & known reconstruction support \\
\texttt{valid\_mask} & $(256,256)$ & usable measurement pairs \\
\texttt{metadata} & scalar/vector & geometry, timing, and provenance \\
\bottomrule
\end{tabular}
\end{center}

Zarr or HDF5 will be used for chunked, compressed, and resumable storage. At
512 temporal samples, one float16 trace cube requires approximately
$67\,\mathrm{MB}$ before compression. One thousand examples therefore require
approximately $67\,\mathrm{GB}$ for traces alone. Raw traces will be retained
as archival data, while compact frequency-domain features will be used during
training.

The planned dataset progression is:
\begin{enumerate}
  \item 16 examples for end-to-end pipeline testing;
  \item 200 training, 25 validation, and 25 test examples for a pilot study;
  \item at least 1000 diverse training examples for the main model.
\end{enumerate}

All slices from one patient must remain in the same split. The patient split
will occur before slice selection or augmentation to prevent information
leakage.

\section{Pulse and Waveform Representation}

Passing all $T$ time samples to a two-dimensional TFNO would be expensive and
would make storage and training unnecessarily difficult. Instead, the pulse
and received signals will be transformed into a stabilized complex frequency
response. For selected frequencies $f_k$,
\begin{equation}
  H_{ij}(f_k)=
  \frac{D_{ij}(f_k)S^*(f_k)}{|S(f_k)|^2+\epsilon},
\end{equation}
where $S$ and $D_{ij}$ are the Fourier transforms of the transmitted and
received signals, respectively. The regularization $\epsilon$ prevents
division by low pulse energy.

The initial feature configuration will use 16 frequencies in the transmitted
band, nominally $100$--$250\,\mathrm{kHz}$. The TFNO input will include:
\begin{itemize}
  \item real and imaginary parts of $H_{ij}(f_k)$: 32 channels;
  \item real and imaginary parts of the normalized pulse spectrum, broadcast
        over the transmitter-receiver grid: 32 channels;
  \item $\sin\theta_{\mathrm{tx}}$, $\cos\theta_{\mathrm{tx}}$,
        $\sin\theta_{\mathrm{rx}}$, and $\cos\theta_{\mathrm{rx}}$;
  \item normalized transmitter-receiver angular separation;
  \item the measurement validity mask.
\end{itemize}

The resulting initial tensor has approximately 70 channels:
\begin{equation}
  X\in\mathbb{R}^{B\times70\times256\times256}.
\end{equation}
This remains a full-waveform approach because complex amplitude and phase are
retained at multiple frequencies; it is not limited to a picked arrival time.

\section{TFNO Reconstruction Architecture}

The model will contain three modules:
\begin{enumerate}
  \item a deterministic pulse-and-trace feature builder;
  \item a TFNO backbone operating on the periodic transmitter-receiver grid;
  \item a learned Cartesian decoder producing the physical image grid.
\end{enumerate}

The pilot TFNO configuration will use:
\begin{itemize}
  \item input grid: $256\times256$ transmitter-receiver samples;
  \item Fourier modes: $32\times32$, with $48\times48$ tested later;
  \item hidden channels: 32, with 64 tested later;
  \item four Fourier layers;
  \item GELU activations;
  \item Tucker-factorized spectral weights;
  \item tensor rank between 0.1 and 0.2.
\end{itemize}

The TFNO produces a latent $256\times256$ field. A residual convolutional
decoder then changes from the angular transmitter-receiver representation to
the Cartesian image representation and performs learned $2\times$ upsampling
to $512\times512$. Keeping this decoder explicit avoids relying on
library-specific super-resolution behavior.

\section{Output Parameterization}

The network will predict slowness-squared rather than speed directly:
\begin{equation}
  m(x,y)=\frac{1}{c(x,y)^2}.
\end{equation}
More specifically, it will predict contrast relative to the water background,
\begin{equation}
  \delta m(x,y)=m(x,y)-m_{\mathrm{water}}.
\end{equation}
The final speed estimate is
\begin{equation}
  \widehat{c}(x,y)=\frac{1}{\sqrt{\widehat{m}(x,y)}}.
\end{equation}
This parameterization is consistent with the existing adjoint inversion,
relates naturally to travel time, and allows the output to be bounded to the
speed range represented in the training data. Pixels outside the known ring
will be set to coupling-water speed rather than predicted.

\section{Training Objective}

The initial loss will combine physical-value, boundary, and spectral terms:
\begin{equation}
  \mathcal{L}=
  \lambda_m\mathcal{L}_{\mathrm{Huber}}(m,\widehat{m})
  +\lambda_c\lVert c-\widehat{c}\rVert_1
  +\lambda_g\mathcal{L}_{\mathrm{grad}}
  +\lambda_f\mathcal{L}_{\mathrm{Fourier}}.
\end{equation}
All terms will be masked to the interior of the ring. Tissue-balanced weights
will prevent the large water and soft-tissue regions from overwhelming bone,
fat, and small anatomical structures. The gradient term will use Sobel
derivatives, while the Fourier term will compare spatial spectra at low and
mid frequencies.

Training will begin with AdamW, automatic mixed precision, batch size one or
two, and gradient accumulation. Tucker factorization reduces parameter memory,
but activations on a $256\times256$ grid will still dominate GPU use.

\section{Physically Valid Augmentation}

The training pipeline will use transformations that respect the ring
geometry:
\begin{itemize}
  \item joint cyclic shifts of transmitter and receiver indices with the
        corresponding target-image rotation;
  \item ring reflections with the correct element-index reversal;
  \item additive electronic noise;
  \item receiver gain variation;
  \item timing jitter;
  \item random dead channels accompanied by the validity mask;
  \item moderate pulse-frequency, bandwidth, phase, and amplitude variation.
\end{itemize}
Whenever the pulse is varied, the exact pulse remains an input to the model.

\section{Evaluation}

The test set will contain only patients absent from training and validation.
Reported metrics will include:
\begin{itemize}
  \item mean absolute error in $\mathrm{m/s}$ inside the ring;
  \item relative $L_2$ error;
  \item structural similarity index;
  \item spatial-gradient error;
  \item tissue-specific error for bone, fat, muscle, and visceral tissue;
  \item robustness to noise, timing jitter, gain error, and missing channels;
  \item inference time, parameter count, and peak GPU memory.
\end{itemize}

The principal comparison will be against the existing nonlinear adjoint FWI.
A homogeneous-water reconstruction and a similarly sized U-Net will serve as
additional baselines. Ablations will compare full complex frequency features,
time-of-flight features, and removal of pulse conditioning.

\section{Proposed Software Structure}

\begin{verbatim}
fdtd2d/
    acquisition.py          Full 256-by-256 ring acquisition

ring_tfno/
    config.py               Geometry and model configuration
    storage.py              Zarr/HDF5 sample format
    features.py             Pulse/trace FFT and frequency response
    dataset.py              PyTorch dataset and normalization
    model.py                TFNO and Cartesian decoder
    losses.py               Physical and image-space losses
    metrics.py              Reconstruction metrics
    checkpoint.py           Reproducible checkpoints

generate_tfno_dataset.py    Generate full-ring examples
train_tfno.py               Train and validate
evaluate_tfno.py            Held-out evaluation and comparisons
infer_tfno.py               Reconstruct one acquisition
\end{verbatim}

The implementation will use PyTorch and the official
\texttt{neuraloperator} package, whose TFNO implementation supports Tucker
factorization and factorized spectral contractions~\cite{kossaifi2024}.
Dependency versions and the complete training configuration will be saved with
every checkpoint.

\section{Implementation Sequence}

\begin{enumerate}
  \item Freeze and test the common $512\times512$ geometry.
  \item Implement one complete 256-transmitter acquisition.
  \item Verify dimensions, timing, reciprocity, and signal normalization.
  \item Define resumable chunked storage and metadata validation.
  \item Generate the 16-example smoke-test dataset.
  \item Implement frequency-response feature extraction.
  \item Implement the TFNO backbone and Cartesian decoder.
  \item Overfit two examples as a model and loss correctness test.
  \item Train and evaluate the 200-example pilot dataset.
  \item Compare the pilot against adjoint FWI and U-Net.
  \item Scale dataset generation only after the pilot demonstrates learning.
\end{enumerate}

\section{Acceptance Criteria}

The first acquisition milestone is complete when one sample:
\begin{itemize}
  \item contains a finite $(256,256,T)$ trace tensor;
  \item uses the same geometry as every future sample;
  \item passes reciprocity and metadata checks;
  \item can be interrupted and resumed without recomputing finished shots;
  \item produces deterministic frequency features of the expected shape.
\end{itemize}

The pilot-model milestone is complete when the model can deliberately overfit
two examples, trains stably on the pilot split, reconstructs held-out anatomy
better than homogeneous water, and is evaluated using the same test cases as
the adjoint FWI baseline.

\begin{thebibliography}{9}

\bibitem{dai2023}
H.~Dai, M.~Penwarden, R.~M. Kirby, and S.~Joshi,
``Neural Operator Learning for Ultrasound Tomography Inversion,''
\emph{Medical Imaging with Deep Learning}, 2023.
arXiv:2304.03297.
\url{https://arxiv.org/abs/2304.03297}

\bibitem{li2021}
Z.~Li, N.~Kovachki, K.~Azizzadenesheli, B.~Liu, K.~Bhattacharya,
A.~Stuart, and A.~Anandkumar,
``Fourier Neural Operator for Parametric Partial Differential Equations,''
\emph{International Conference on Learning Representations}, 2021.
\url{https://openreview.net/forum?id=c8P9NQVtmnO}

\bibitem{kovachki2023}
N.~Kovachki, Z.~Li, B.~Liu, K.~Azizzadenesheli, K.~Bhattacharya,
A.~Stuart, and A.~Anandkumar,
``Neural Operator: Learning Maps Between Function Spaces with Applications
to PDEs,'' \emph{Journal of Machine Learning Research}, vol.~24, no.~89,
pp.~1--97, 2023.
\url{https://www.jmlr.org/papers/v24/21-1524.html}

\bibitem{kossaifi2024}
J.~Kossaifi, N.~Kovachki, Z.~Li, D.~Pitt, M.~Liu-Schiaffini,
R.~J. George, B.~Bonev, K.~Azizzadenesheli, J.~Berner, and A.~Anandkumar,
``A Library for Learning Neural Operators,'' 2024.
arXiv:2412.10354.
\url{https://arxiv.org/abs/2412.10354}

\bibitem{tran2021}
A.~Tran, A.~Mathews, L.~Xie, and C.~S. Ong,
``Factorized Fourier Neural Operators,'' 2021.
arXiv:2111.13802.
\url{https://arxiv.org/abs/2111.13802}

\end{thebibliography}

\end{document}
